\documentclass[11pt,a4paper]{article}

% Page and text formatting follow the supplied literature review.
\usepackage[a4paper,top=25.6mm,right=25.4mm,bottom=12.5mm,left=25.4mm]{geometry}
\usepackage[T1]{fontenc}
\usepackage[utf8]{inputenc}
\usepackage{newtxtext,newtxmath}
\usepackage{setspace}
\usepackage{microtype}
\usepackage[hidelinks]{hyperref}
\usepackage{enumitem}

\setstretch{1.15}
\setlength{\parindent}{0pt}
\setlength{\parskip}{0.75em}
\setlist[itemize]{leftmargin=1.7em,itemsep=0.3em,topsep=0.25em}
\renewcommand{\contentsname}{Table of Contents}

\begin{document}

\begin{titlepage}
    \centering
    \vspace*{1.2cm}

    {\fontsize{28}{34}\selectfont An XR System for Concentration Training Under Exam Conditions\par}
    \vspace{0.8cm}
    {\fontsize{16}{20}\selectfont\bfseries Mid-Year Progress Report\par}

    \vspace{1.8cm}
    {\fontsize{16}{20}\selectfont\bfseries Perry Xie\par}

    \vspace{1.4cm}
    {\fontsize{14}{18}\selectfont
    Department of Electrical, Computer, and Software Engineering\par
    The University of Auckland, Auckland, New Zealand\par}

    \vfill
    {\fontsize{16}{20}\selectfont\bfseries July 2026\par}

    \vspace{1.4cm}
    {\fontsize{16}{20}\selectfont
    Supervisor: Yun Suen Pai\par
    Project Mentor: Yuzhen Gao\par
    Project Partner: Weyman Wong\par}

    \vfill
    \begin{minipage}{0.9\textwidth}
        \centering\fontsize{10}{13}\selectfont\itshape\sffamily
        Disclaimer -- This report has been submitted as partial fulfilment of the
        requirements for the BE(Hons) Part IV research project.
    \end{minipage}
\end{titlepage}

\clearpage
\section*{Declaration of Originality}
This report is my own unaided work and was not copied from, nor written in
collaboration with, any other person except where collaboration and sources are
explicitly acknowledged.

\vspace{2em}
Name: Perry Xie

\clearpage
\tableofcontents

\clearpage
\section{Introduction}

Examinations are a central form of academic assessment, yet the conditions in
which they are completed can make sustained concentration difficult. Students
must manage time pressure, cognitive load, and stress while resisting both
internal distractions, such as mind-wandering, and external distractions, such
as surrounding movement and noise. When attention shifts away from the task,
students may not immediately recognise the lapse or successfully redirect their
focus. This can reduce task performance and limit their ability to demonstrate
their knowledge under exam conditions \cite{smallwood2015,mooneyham2013}.

Virtual reality (VR) offers a controlled and repeatable way to recreate these
conditions for attention training. An immersive exam-like environment can
present realistic auditory and visual distractions while allowing their timing,
type, and intensity to be manipulated consistently. Physiological and
behavioural measures can also provide real-time evidence of changes in a user's
attentional state. In particular, eye-gaze behaviour can indicate whether visual
attention has moved away from the task, while heart-rate variability can provide
complementary information about stress and self-regulation. Combining these
signals creates an opportunity to detect probable attention lapses as they occur,
rather than relying only on retrospective self-report.

This project investigates an XR concentration-training system designed for
exam-like conditions. The system combines a VR examination environment,
controlled visual and auditory distractions, and real-time physiological
monitoring. It applies the Just-in-Time Adaptive Intervention (JITAI) framework
to determine when support is required and to deliver a contextually appropriate
intervention intended to help the user re-engage with the task
\cite{nahumshani2018}. The central design challenge is to make this feedback
timely enough to support attention recovery while keeping it subtle enough that
the intervention does not become an additional distraction.

The aim of the project is to design and evaluate a VR-based biofeedback system
that helps students regulate attention and maintain focus under simulated exam
conditions. The work focuses on three connected capabilities: creating a
realistic distraction-rich VR environment; using eye-gaze and heart-rate data to
identify changes associated with attention and stress; and implementing adaptive
feedback that responds to the user's current state. Together, these components
form a closed-loop training system in which attention is monitored, an
intervention is selected when required, and the user's subsequent response can
inform continued adaptation.

This mid-year progress report presents the development completed to date and the
work planned for the remainder of the project. It describes the system design,
the implementation of the VR environment and distraction mechanisms, the
integration of physiological data, and the development of the JITAI feedback
logic. It also identifies current limitations, technical risks, and the testing
required before the system can be evaluated against the project's research
questions.

\begin{thebibliography}{9}

\bibitem{smallwood2015}
J. Smallwood and J. W. Schooler,
``The science of mind wandering: Empirically navigating the stream of
consciousness,''
\emph{Annual Review of Psychology}, vol. 66, pp. 487--518, 2015.

\bibitem{mooneyham2013}
B. W. Mooneyham and J. W. Schooler,
``The costs and benefits of mind-wandering: A review,''
\emph{Canadian Journal of Experimental Psychology}, vol. 67, no. 1,
pp. 11--18, 2013.

\bibitem{nahumshani2018}
I. Nahum-Shani, S. N. Smith, B. J. Spring, L. M. Collins, K. Witkiewitz,
A. Tewari, and S. A. Murphy,
``Just-in-time adaptive interventions (JITAIs) in mobile health: Key
components and design principles for ongoing health behavior support,''
\emph{Annals of Behavioral Medicine}, vol. 52, no. 6, pp. 446--462, 2018.

\end{thebibliography}

\end{document}
